% !TEX root = ../libretto.tex
% Sezione: Ingresso
\momento{Ingresso}
\begin{songs}{titoli}

% Fonte: LibrettoAccordi.pdf (Canti Coro), A4.
% Nell'elenco: Ingresso, Ordinario.
\beginsong{Acclamate al Signore}[by={M. Frisina}]
\beginchorus
Acclamate al Signore,
voi tutti della terra,
e servitelo con gioia,
andate a lui con esultanza.
Acclamate voi tutti
al Signore.
\endchorus
\beginverse
Riconoscete che il Signore, che il Signore è Dio,
Egli ci ha fatti, siamo suoi, suo popolo e gregge del suo pascolo.
\endverse
\beginverse
Entrate nelle sue porte con degli inni di grazia,
suoi atri nella lode, benedite, lodate il suo nome.
\endverse
\beginverse
Poiché buono è il Signore, eterna la sua misericordia,
la sua fedeltà si estende sopra ogni generazione.
\endverse
\endsong

% Fonte: LibrettoAccordi.pdf (Canti Coro), J42.
% Nell'elenco: Ingresso.
\beginsong{Beati voi}[by={F. Buttazzo}]
\beginchorus
Beati voi! Beati voi! Beati voi! Beati voi!
\endchorus
\beginverse
Se sarete poveri nel cuore, beati voi:
sarà vostro il Regno di Dio Padre.
Se sarete voi che piangerete, beati voi,
perché un giorno vi consolerò.
\endverse
\beginverse
Se sarete miti verso tutti, beati voi:
erediterete tutto il mondo.
Quando avrete fame di giustizia, beati voi,
perché un giorno Io vi sazierò.
\endverse
\beginverse
Se sarete misericordiosi, beati voi:
la misericordia troverete.
Se sarete puri dentro il cuore, beati voi,
perché voi vedrete il Padre mio.
\endverse
\beginverse
Se lavorerete per la pace, beati voi:
chiameranno voi «figli di Dio».
Se per causa mia voi soffrirete, beati voi,
sarà grande in voi la santità.
\endverse
\endsong

% Fonte: LibrettoAccordi.pdf (Canti Coro), A9.
% Nell'elenco: Ingresso.
\beginsong{Cantate al Signore}[by={Gen Rosso}]
\beginchorus
Alleluja! Alleluja! Alleluja! Alleluja! \rep{2}
\endchorus
\beginverse
Cantate al Signore un cantico nuovo, tutta la terra canti al Signore,
per tutta la terra si è accesa una luce: uomini nuovi cantano in cor
un cantico nuovo di gioia infinita, un canto d'amore a Dio fra noi.
Alleluja!
\endverse
\beginverse
Un coro di voci s'innalza al Signore, re della vita, luce del mondo;
discende dal cielo un fuoco d'amore: il Paradiso canta con noi
un cantico nuovo di gioia infinita, un canto d'amore a Dio fra noi.
Alleluja!
\endverse
\endsong

% Fonte: testo fornito dall'utente (tolti gli accordi). Non presente nell'elenco.
% Nel foglio il ritornello è scritto due volte di seguito; qui una volta sola. Finale: ultimo verso ripetuto.
\beginsong{Cantiamo al Signore con gioia}[by={F. Volpintesta}]
\beginchorus
Cantiamo al Signore con gioia,
cantiamo per amor:
un inno di lode la vita sarà,
un canto d'amore per lui.
\endchorus
\beginverse
In ogni momento, un dono,
la vita per sempre sarà,
e nasce la gioia che viene da te:
la gioia che mai finirà.
\endverse
\beginverse
La vita sarà solo amore
se tu resti ancora con noi:
la tua presenza la luce darà,
la luce che viene da te.
\endverse
\beginchorus
Cantiamo al Signore con gioia,
cantiamo per amor:
un inno di lode la vita sarà,
un canto d'amore per lui. \rep{2}
\endchorus
\endsong

% Fonte: LibrettoAccordi.pdf (Canti Coro), A2.
% Nell'elenco: Ingresso, Ordinario.
\beginsong{Cantiamo con gioia}[by={RnS}]
\beginchorus
Cantiamo con gioia al nostro Dio,
gridiamo con forza al nostro Re!
Le mani battiamo, insieme danziamo, cantiamo per la felicità!
\endchorus
\beginverse
Rialza chiunque è caduto, rafforza chi è debole tra noi.
Chi è triste ritrova in Lui la gioia.
Su canta insieme a noi e grida più che puoi
che Gesù Cristo è il Signor!
\endverse
\beginverse
Eterna è la sua misericordia, fedele per sempre è il suo amor.
Con forza ha sconfitto il nemico.
Su canta insieme a noi e grida più che puoi
che Gesù Cristo è il Signor!
\endverse
\endsong

% Testo fornito dall'utente (fonte: genverde.it/en/?p=34261).
% Versione italiana di "You Call Us by Name".
% Struttura cantata: Rit. - Strofa 1 - Rit. - Strofa 2 - Rit. + finale.
% Nell'elenco: Ingresso, Ordinario.
\beginsong{Chiamati per nome}[by={Gen Verde (J. Bai, N. L. Uelmen, S. G. McAllister)}]
\beginchorus
Veniamo da te,
chiamati per nome.
Che festa, Signore, tu cammini con noi.
Ci parli di te,
per noi spezzi il pane,
ti riconosciamo e il cuore arde, sei tu.
E noi tuo popolo
siamo qui.
\endchorus
\beginverse
Siamo come terra ed argilla
e la tua parola ci plasmerà;
brace pronta per la scintilla
e il tuo spirito soffierà,
c'infiammerà.
\endverse
\beginverse
Siamo come semi nel solco,
come vigna che il suo frutto darà;
grano del Signore risorto,
la tua messe che fiorirà d'eternità.
\endverse
\beginverse*
\textit{Finale (dopo l'ultimo ritornello):}
E noi tuo popolo
siamo qui, siamo qui.
\endverse
\endsong

% Fonte: LibrettoAccordi.pdf (Canti Coro), A19.
% Nell'elenco: Ingresso, Matrimonio (Baggio-Buttazzo). CD: Dall'amore di Dio (2010).
\beginsong{Dall'amore di Dio}[by={F. Buttazzo (CD Dall'amore di Dio)}]
\beginverse
Nasce dall'amore di Dio il miracolo di questa vita.
Vive per l'amore di Dio ogni angolo dell'universo.
Ma nel cuore d'ogni uomo risplende
il dono più divino: la vita sua più vera.
\endverse
\beginchorus
Creati come te, amati noi da te,
voliamo più su, più in alto di ogni cosa.
Creati come te, chiamati noi da te,
doniamo la vita con te.
Creati come te, amati noi da te,
voliamo più su, più in alto di ogni cosa.
Creati come te, chiamati noi da te,
amiamo per l'eternità.
\endchorus
\beginverse
Nasce dentro il cuore di Dio il sentiero della nostra vita.
Viene dall'amore di Dio la bellezza della nostra storia.
Noi, chiamati a donarci la vita,
seguiamo in due una strada, vediamo in lui la meta.
\endverse
\endsong

% Fonte: LibrettoAccordi.pdf (Canti Coro), A5.
% Nell'elenco: Ingresso, Ordinario.
\beginsong{Dall'aurora al tramonto}[by={M. Balduzzi (CD Verbum Panis)}]
\beginchorus
Dall'aurora io cerco te, fino al tramonto ti chiamo
ha sete solo di te l'anima mia come terra deserta.
\endchorus
\beginverse
Non mi fermerò un solo istante sempre canterò la tua lode
perché sei il mio Dio, il mio riparo
mi proteggerai all'ombra delle tue ali.
\endverse
\beginverse
Non mi fermerò un solo istante io racconterò le tue opere
perché sei il mio Dio, unico bene
nulla mai potrà la notte contro di me.
\endverse
\endsong

% Fonte: LibrettoAccordi.pdf (Canti Coro), A1.
% Nell'elenco: Ingresso, Ordinario.
\beginsong{È bello lodarti}[by={Gen Verde}]
\beginchorus
È bello cantare il tuo amore, è bello lodare il tuo nome.
È bello cantare il tuo amore, è bello lodarti Signore,
è bello cantare a te. \rep{2}
\endchorus
\beginverse
Tu che sei l'Amore infinito, che neppure il cielo può contenere,
ti sei fatto uomo, tu sei venuto qui
ad abitare in mezzo a noi, allora
\endverse
\beginverse
Tu, che conti tutte le stelle e le chiami una ad una per nome,
da mille sentieri ci hai radunati qui,
ci hai chiamati figli tuoi, allora
\endverse
\endsong

% Fonte: LibrettoAccordi.pdf (Canti Coro), O7.
% Nell'elenco: Ingresso.
\beginsong{È giunta l'ora}[by={L. Scaglianti}]
\beginverse
È giunta l'ora, Padre, per me: ai miei amici ho detto che
questa è la vita: conoscere te e il Figlio tuo, Cristo Gesù.
\endverse
\beginverse
Erano tuoi, li hai dati a me ed ora sanno che torno a te.
Hanno creduto: conservali tu nel tuo amore, nell'unità.
\endverse
\beginverse
Tu mi hai mandato ai figli tuoi: la tua parola è verità.
E il loro cuore sia pieno di gioia: la gioia vera viene da te.
\endverse
\beginverse
Io sono in loro e tu in me: che sian perfetti nell'unità
e il mondo creda che tu mi hai mandato, li hai amati come ami me.
\endverse
\endsong

% Fonte: LibrettoAccordi.pdf (Canti Coro), K12.
% Nell'elenco: Ingresso, Ordinario.
\beginsong{È la gioia che fa cantare}[by={RnS}]
\beginchorus
È la gioia che fa cantare, celebrando il Signore
il suo spirito oggi canta in me.
È la gioia che fa cantare, celebrando il Signore
il suo spirito oggi canta in me.
\endchorus
\beginverse
Io canto la gloria tua, perchè hai vinto la morte
Mia potente salvezza, mia forza sei tu.
\endverse
\beginverse
Tu raduni il tuo popolo e sconfiggi le tenebre.
Il tuo esercito siamo noi: «Vittoria di Dio».
\endverse
\endsong

% Fonte: LibrettoAccordi.pdf (Canti Coro), J28.
% Nell'elenco: Ingresso, Ordinario.
\beginsong{Eccomi\soloindice{ (Gen Verde)}}[by={Gen Verde}]
\beginchorus
Eccomi Gesù anche oggi davanti a te,
tutto rinnovato, proprio come Tu mi vuoi,
io sarò la risposta al tuo perché
un frutto degno del tuo abbandono, eccomi. \rep{2}
\endchorus
\beginverse
Un patto ormai ci stringe tutti qui di fronte a te
per dichiararti il nostro amore esclusivo
e corrispondere ai doni del tuo amore
perché sei stato abbandonato, siamo morti e risorti con te.
\endverse
\beginverse
Schierati in prima fila nella divina avventura,
per far conoscere in questo mondo l'amore,
un solo cuore, un solo coro che
ripete una preghiera con la vita dice a Te.
\endverse
\endsong

% Fonte: "99 - Raccolte e libretti accordi/frisina-marco-non-temere-fascicolo-completo_compress.pdf" (pag. 15). Salmo 39.
% Nell'elenco: Ingresso, Ordinario. CD: Non temere, Signore è il suo nome.
\beginsong{Eccomi\soloindice{ (M. Frisina)}}[by={M. Frisina (CD Non temere, Signore è il suo nome)}]
\beginchorus
Eccomi, eccomi,
Signore, io vengo.
Eccomi, eccomi!
Si compia in me la tua volontà.
\endchorus
\beginverse
Nel mio Signore ho sperato
e su di me s'è chinato,
ha dato ascolto al mio grido,
m'ha liberato dalla morte.
\endverse
\beginverse
I miei piedi ha reso saldi,
sicuri ha reso i miei passi.
Ha messo sulla mia bocca
un nuovo canto di lode.
\endverse
\beginverse
Il sacrificio non gradisci,
ma m'hai aperto l'orecchio,
non hai voluto olocausti,
allora ho detto: «Io vengo!».
\endverse
\beginverse
Sul tuo libro di me è scritto:
si compia il tuo volere.
Questo, mio Dio, desidero,
la tua legge è nel mio cuore.
\endverse
\beginverse
La tua giustizia ho proclamato,
non tengo chiuse le labbra.
Non rifiutarmi, Signore,
la tua misericordia.
\endverse
\endsong

% Fonte: LibrettoAccordi.pdf (Canti Coro), A8.
% Nell'elenco: Ingresso, Ordinario.
\beginsong{Giorno di festa}[by={C. Bertoni, F. Righetto}]
\beginchorus
È giorno di festa torniamo all'Amore
è giorno di gioia, è Dio fra noi.
Nella tua casa ogni attimo è pace, incontrare Te,
è dono e vita per me.
\endchorus
\beginverse
Come un fiume torna al mare, come un fiore guarda al cielo,
i nostri occhi, le nostre mani,
si alzano a Te, ritrovano Te, che sei Amore.
\endverse
\beginverse
Riconoscerti presente, Tu parola, Tu in quel pane.
Ardente fuoco, brucia nei cuori
che cercano Te, ritrovano Te, che sei Amore.
\endverse
\endsong

% Fonte: LibrettoAccordi.pdf (Canti Coro), D29.
% Nell'elenco: Ingresso, Ordinario.
\beginsong{I cieli narrano}[by={M. Frisina}]
\beginchorus
I cieli narrano la gloria di Dio
e il firmamento annunzia l'opera Sua.
Alleluia, alleluia, alleluia, alleluia.
\endchorus
\beginverse
Il giorno al giorno ne affida il messaggio,
la notte alla notte ne trasmette notizia,
non è linguaggio, non sono parole
di cui non si oda il suono.
\endverse
\beginverse
Là pose una tenda per il sole che sorge,
è come uno sposo dalla stanza nuziale,
esulta come un prode che corre
con gioia la sua strada.
\endverse
\beginverse
Lui sorge dall'ultimo estremo del cielo
e la sua corsa l'altro estremo raggiunge.
Nessuna delle creature potrà
mai sottrarsi al suo calore.
\endverse
\beginverse
La legge di Dio rinfranca l'anima mia,
la testimonianza del Signore è verace,
gioisce il cuore ai suoi giusti precetti
che danno la luce agli occhi.
\endverse
\endsong

% Fonte: LibrettoAccordi.pdf (Canti Coro), A10.
% Nell'elenco: Ingresso.
\beginsong{Lodate il Signore dai cieli}[by={M. Frisina}]
\beginchorus
Lodate il signore dai cieli, nell'alto dei cieli lodatelo.
Lodatelo, voi suoi angeli, lodatelo, voi sue schiere.
\endchorus
\beginverse
Lodatelo, sole e luna, lodatelo, fulgide stelle.
Lodatelo, cieli dei cieli, voi acque al di sopra dei cieli.
\endverse
\beginverse
Lodate il Signore, creature, voi mostri marini e abissi,
il fuoco, la neve, la nebbia, il vento che a lui obbedisce.
\endverse
\beginverse
Voi monti e voi tutte colline, voi alberi e tutti voi cedri,
voi fiere e tutte le bestie, voi rettili e uccelli alati.
\endverse
\beginverse
I re della terra e i popoli, i giudici e i governanti,
i giovani e le fanciulle, i vecchi insieme ai bambini.
\endverse
\beginverse
Perché il suo nome è sublime, la gloria sua splende sul mondo.
È lode per i suoi fedeli, il popolo che egli ama.
\endverse
\endsong

% Fonte: testo fornito dall'utente. Non presente nell'elenco.
\beginsong{Open the eyes of my heart}[by={M. W. Smith}]
\beginverse*
Open the eyes of my heart, Lord,
open the eyes of my heart:
I want to see you,
I want to see you.
\endverse
\beginverse*
To see you high and lifted up,
shinin' in the light of your glory.
Pour out your power and love
as we sing holy, holy, holy.
\endverse
\beginverse*
Holy, holy, holy,
we cry holy, holy, holy,
you are holy, holy, holy.
I want to see you.
\endverse
\endsong

% Fonte: LibrettoAccordi.pdf (Canti Coro), J33.
% Nell'elenco: senza tipologia (nel libretto accordi tra i canti di comunione). L'utente lo mette in Ingresso.
\beginsong{Ora è tempo di gioia}[by={Gen Rosso}]
\beginverse
L'eco torna d'antiche valli, la sua voce non porta più
ricordo di sommesse lacrime di esili in terre lontane.
\endverse
\beginchorus
Ora è tempo di gioia, non ve ne accorgete?
Ecco faccio una cosa nuova,
nel deserto una strada aprirò.
\endchorus
\beginverse
Come l'onda che sulla sabbia copre le orme e poi passa e va,
così nel tempo si cancellano le ombre scure del lungo inverno.
\endverse
\beginverse
Fra i sentieri dei boschi il vento con i rami ricomporrà
nuove armonie che trasformano i lamenti in canti di festa.
\endverse
\endsong

% Fonte: LibrettoAccordi.pdf (Canti Coro), J40 (li' tra i canti di comunione).
% Nell'elenco: Ingresso, Ordinario. CD: Venti dello Spirito (1988-2007).
\beginsong{Popoli tutti, acclamate}[by={RnS (CD Venti dello Spirito)}]
\beginverse
Mio Dio, Signore, nulla è pari a Te.
Ora e per sempre, voglio lodare il tuo grande amor per noi.
Mia roccia Tu sei, pace e conforto mi dai,
con tutto il cuore e le mie forze, sempre io ti adorerò.
\endverse
\beginchorus
Popoli tutti, acclamate al Signore,
gloria e potenza cantiamo al Re,
mari e monti si prostrino a Te, al tuo nome, o Signore.
Canto di gioia per quello che fai,
per sempre, Signore, con Te resterò,
non c'è promessa, non c'è fedeltà che in Te.
\endchorus
\endsong

% Fonte: LibrettoAccordi.pdf (Canti Coro), A21.
% Nell'elenco: Ingresso, Ordinario.
\beginsong{Quale gioia è star con te}[by={RnS}]
\beginverse
Ogni volta che ti cerco, ogni volta che t'invoco,
sempre mi accogli, Signor.
Grandi sono i tuoi prodigi, tu sei buono verso tutti,
santo tu regni tra noi.
\endverse
\beginchorus
Quale gioia è star con te, Gesù, vivo e vicino,
bello è dar lode a Te, Tu sei il Signor.
Quale dono è aver creduto in Te, che non mi abbandoni,
io per sempre abiterò la tua casa, mio Re.
\endchorus
\beginverse
Hai guarito il mio dolore, hai cambiato questo cuore,
oggi rinasco, Signor.
Grandi sono i tuoi prodigi, tu sei buono verso tutti,
santo tu regni tra noi.
\endverse
\beginverse
Hai salvato la mia vita, hai aperto la mia bocca,
canto per te, mio Signor.
Grandi sono i tuoi prodigi, tu sei buono verso tutti,
santo tu regni tra noi.
\endverse
\beginverse*
Ti loderò, ti adorerò, ti canterò,
noi ti canteremo.
\endverse
\endsong

% Fonte: LibrettoAccordi.pdf (Canti Coro), J19; c'e' anche
% "1 Canti di Ingresso/Una goccia nel mare (M. Balduzzi).pdf".
% Nell'elenco: Ingresso, Ordinario. CD: Xaipe.
\beginsong{Una goccia del mare}[by={M. Balduzzi (CD Xaipe)}]
\beginverse
Tu che per nome chiami tutte le stelle,
tu che raccogli ogni nostro segreto sospiro,
non mi abbandonare,
non mi lasciare nell'ombra del cuore,
e da te io verrò.
\endverse
\beginchorus
Una goccia del mare ti chiedo,
che disseti come rugiada,
dammi solo una goccia del mare
dell'infinita sapienza tua.
\endchorus
\beginverse
Tu che hai creato il calice dell'universo,
tu che per gli uomini hai preparato una casa,
non mi abbandonare,
non mi lasciare nell'ombra del cuore
e da te io verrò.
\endverse
\endsong

% Fonte: testo fornito dall'utente. Non presente nell'elenco. Autore non indicato.
\beginsong{Vida en abundancia}
\beginverse
Los lirios del campo
y las aves del cielo
no se preocupan
porque están en mis manos.
Tené confianza en mí,
acá estoy junto a vos.
Amá lo que sos y tus circunstancias,
estoy con vos, con tu cruz en mi espalda.
Todo terminará bien,
yo hago nuevas todas las cosas.
\endverse
\beginchorus
Yo vengo a traerte vida,
vida en abundancia, en abundancia.
Yo soy el camino, la verdad y la vida,
vida en abundancia, en abundancia.
\endchorus
\beginverse
No hice al hombre para que esté solo:
caminen juntos, como hermanos,
sopórtense mutuamente,
ámense unos a otros.
La felicidad de la vida eterna
empieza conmigo en la tierra.
Sentite vivo:
la fiesta del reino comienza acá.
\endverse
\endsong

% Fonte: LibrettoAccordi.pdf (Canti Coro), D8.
% Nell'elenco: Ingresso, Matrimonio.
\beginsong{Vieni dal Libano}[by={K. Arguello}]
\beginchorus
Vieni dal Libano, mia sposa, vieni dal Libano, vieni!
Avrai per corona le vette dei monti, le alte cime dell'Hermon.
Tu m'hai ferito, ferito il cuore, o sorella mia sposa.
Vieni dal Libano, mia sposa, vieni dal Libano, vieni.
\endchorus
\beginverse*
Cercai l'amore dell'anima mia, lo cercai senza trovarlo.
Trovai l'amore dell'anima mia, l'ho abbracciato e non lo lascerò mai!
\endverse
\beginverse
Alzati in fretta, mia diletta, vieni, colomba, vieni.
L'inverno ormai è già passato, il canto della tortora si ode.
I fiori son tornati sulla terra, il grande sole è venuto.
Alzati in fretta, mia diletta, vieni, colomba, vieni.
\endverse
\beginverse
Come sigillo sul tuo cuore, come sigillo sul braccio;
che l'amore è forte come la morte e le acque non lo spegneranno.
Dare per esso tutti i beni della casa sarebbe disprezzarlo.
Come sigillo sul tuo cuore, come sigillo sul braccio.
\endverse
\endsong

\end{songs}
